% !TEX program = xelatex
% 공통수학2 · Ⅱ. 집합과 명제 · 1. 집합 · 03 교집합과 합집합 · 2/2차시
\documentclass[11pt,a4paper]{article}
\usepackage{kotex}
\usepackage{iftex}
\ifPDFTeX\else
  \setmainhangulfont{Noto Serif CJK KR}
  \setsanshangulfont{Noto Sans CJK KR}
\fi
\usepackage[margin=2.1cm]{geometry}
\usepackage{parskip}
\usepackage{enumitem}
\usepackage{array}
\usepackage{tabularx}
\usepackage{booktabs}
\usepackage{xcolor}
\usepackage{amssymb}
\usepackage{tikz}
\usepackage[most]{tcolorbox}
\usepackage{fancyhdr}
\usepackage{titlesec}

\definecolor{goldc}{HTML}{B07C22}
\definecolor{goldstrong}{HTML}{8A5F16}
\definecolor{indigoc}{HTML}{2B3B57}
\definecolor{indigosoft}{HTML}{6E82A6}
\definecolor{linec}{HTML}{D6DACB}
\definecolor{surface2}{HTML}{E4E8DC}
\definecolor{notebg}{HTML}{FBF3E3}
\definecolor{inksoft}{HTML}{52604F}

\titleformat{\section}{\normalfont\sffamily\bfseries\large\color{indigoc}}{\thesection}{0.6em}{}
\titlespacing{\section}{0pt}{16pt}{8pt}
\newtcolorbox{ideabox}{enhanced, breakable, colback=white, colframe=goldc, boxrule=0.5pt, leftrule=3pt, arc=2pt, left=10pt, right=10pt, top=8pt, bottom=8pt}
\newtcolorbox{calloutbox}{enhanced, breakable, colback=white, colframe=indigosoft, boxrule=0.5pt, leftrule=3pt, arc=2pt, left=10pt, right=10pt, top=7pt, bottom=7pt}
\newtcolorbox{notebox}{enhanced, breakable, colback=notebg, colframe=goldc, boxrule=0.5pt, leftrule=3pt, arc=2pt, left=10pt, right=10pt, top=7pt, bottom=7pt}
\newtcolorbox{answerbox}{enhanced, breakable, colback=white, colframe=linec, boxrule=0.5pt, arc=2pt, left=10pt, right=10pt, top=7pt, bottom=7pt}
\newcommand{\lessonbanner}[3]{%
  \begin{tcolorbox}[enhanced, colback=indigoc, colframe=indigoc, boxrule=0pt, arc=0pt,
    left=14pt, right=14pt, top=14pt, bottom=10pt, borderline south={2.4pt}{0pt}{goldc}, fontupper=\color{white}]
    {\sffamily\footnotesize\color{white!85!black} #1}\\[4pt]{\sffamily\bfseries\Large #2}\\[3pt]{\sffamily\small\color{white!88!black} #3}
  \end{tcolorbox}}
\pagestyle{fancy}\fancyhf{}\renewcommand{\headrulewidth}{0pt}
\fancyfoot[L]{\footnotesize\color{inksoft} 공통수학2 · 2026학년도 2학기 · 지평선고등학교}
\fancyfoot[R]{\footnotesize\color{inksoft} 교사용 지도안 -- 배포 금지}

\begin{document}
\lessonbanner{공통수학2 · Ⅱ. 집합과 명제 · 1. 집합}{03 교집합과 합집합 --- 2/2차시}{원소의 개수 공식과 연산 법칙 · 교사용 지도안 (2026학년도 2학기)}
\vspace{10pt}

\noindent\begin{tabularx}{\linewidth}{@{}>{\bfseries\color{indigoc}}p{2.6cm}X@{}}
\toprule
대상 & 고1 · 2개 반 · 반당 13명 \\
차시 & 50분 (도입 10 · 전개 30 · 정리 10) \\
성취기준 & 집합의 연산을 수행하고, 벤 다이어그램을 이용하여 나타낼 수 있다. \\
선행 학습 & 21차시 --- 교집합과 합집합의 뜻 \\
\bottomrule
\end{tabularx}

\section{학습 목표 \& Big Idea}
\begin{ideabox}
\textbf{\color{goldstrong}영속적 이해 (Big Idea)}\\
$n(A)+n(B)$는 교집합 부분을 두 번 세므로, 한 번 빼 주어야 정확한 개수가 된다. 또한 집합의 연산에도 수의 연산처럼 `순서를 바꿔도, 묶는 방법을 바꿔도 결과가 같다'는 법칙들이 있다 --- 그러나 모든 법칙이 그대로 옮겨지지는 않는다.
\vspace{4pt}
\small\color{inksoft} 핵심개념: \textbf{중복 제거} \quad\textbullet\quad 핵심개념: \textbf{연산 법칙} \quad\textbullet\quad 일반화: \textbf{비슷해 보이는 구조도 세부 규칙은 다를 수 있다}
\end{ideabox}
\begin{calloutbox}
\textbf{차시 학습 목표}
\begin{itemize}[leftmargin=1.4em, itemsep=2pt, topsep=4pt]
  \item $n(A\cup B)=n(A)+n(B)-n(A\cap B)$ 공식을 이해하고 활용할 수 있다.
  \item 집합의 연산에 대한 교환법칙, 결합법칙, 분배법칙을 벤 다이어그램으로 확인할 수 있다.
\end{itemize}
\end{calloutbox}

\section{질문의 연결}
\begin{tabularx}{\linewidth}{@{}>{\bfseries\color{indigoc}\small}p{2.4cm}X@{}}
사실적 질문 & $n(A)+n(B)$를 그대로 $n(A\cup B)$라고 할 수 없는 이유는 무엇일까? \\[6pt]
개념적 질문 & 수의 덧셈·곱셈 법칙(교환·결합·분배)이 집합의 합집합·교집합에도 그대로 성립하는데, 왜 $(A\cup B)\cap C \ne A\cup(B\cap C)$일까? \\[6pt]
논쟁적 질문 & 비슷해 보이는 두 구조(수의 연산과 집합의 연산)가 어떤 규칙은 공유하고 어떤 규칙은 공유하지 않는다는 것은 무엇을 말해 줄까? \\
\end{tabularx}

\section{수업 흐름}
\begin{tabularx}{\linewidth}{@{}>{\centering\arraybackslash}p{1.6cm}X>{\raggedright\arraybackslash\small}p{3.1cm}@{}}
\toprule
\textbf{단계} & \textbf{활동 \& 발문} & \textbf{ATL / VTR} \\
\midrule
\textcolor{goldstrong}{\textbf{도입}}\newline\footnotesize 10분 &
\textbf{마음공부 (2분)} --- 유무념 대조.\newline
\textbf{생각 열기 (8분)} --- 벤 다이어그램의 구체적인 두 집합에서 $n(A)$, $n(B)$, $n(A\cap B)$, $n(A\cup B)$를 직접 세어 보고, $n(A)+n(B)-n(A\cap B)$와 비교(함께하기). &
ATL: 사고(수량 비교)\newline VTR: Connect-Extend-Challenge \\
\addlinespace
\textcolor{goldstrong}{\textbf{전개}}\newline\footnotesize 30분 &
\textbf{원소의 개수 공식 (10분)} --- $n(A\cup B)=n(A)+n(B)-n(A\cap B)$ 공식화(중복 제거 원리) $\to$ 문제 3.\newline
\textbf{연산 법칙 (20분)} --- 교환법칙(자명) $\to$ 결합법칙(문제 4, 벤 다이어그램 색칠) $\to$ 분배법칙(문제 5) $\to$ 오개념 확인: $(A\cup B)\cap C \ne A\cup(B\cap C)$ 반례 $A=\{1,2\},B=\{2,3\},C=\{1\}$ 함께 계산. &
ATT: 촉진자 역할\newline ATL: 시각화·검증 \\
\addlinespace
\textcolor{goldstrong}{\textbf{정리}}\newline\footnotesize 10분 &
\textbf{개념 연결 질문 (5분)} --- ``분배법칙과 결합법칙, 벤 다이어그램에서 어떻게 구별해서 그릴까?''\newline
\textbf{성찰 (5분)} --- 감사 일기 1문장 + `교집합과 합집합' 소단원 마무리 + 다음 차시 예고(여집합과 차집합). &
마음공부 · 감사 일기 \\
\bottomrule
\end{tabularx}

\begin{calloutbox}
\textbf{지도상의 유의점} --- $(A\cup B)\cap C$와 $A\cup(B\cap C)$는 서로 다른 집합임을 반드시 반례로 확인시킨다($A=\{1,2\}, B=\{2,3\}, C=\{1\}$일 때 전자는 $\{1\}$, 후자는 $\{1,2\}$). 결합법칙과 분배법칙을 혼동하지 않도록 주의한다. 집합의 분배법칙은 수의 분배법칙과 유사하지만, $a+(b\times c)=(a+b)\times(a+c)$는 성립하지 않는 반면 $A\cup(B\cap C)=(A\cup B)\cap(A\cup C)$는 성립한다는 차이를 짚어 준다.
\end{calloutbox}

\section{형성평가 \& 답안}
\begin{answerbox}
\textbf{문제 3} \quad $n(A)=8$, $n(B)=10$, $n(A\cup B)=16$일 때 $n(A\cap B)$\\
\textbf{\color{goldstrong}답} \; $n(A\cap B)=8+10-16=2$
\end{answerbox}
\begin{answerbox}
\textbf{집합의 연산 법칙 (정리)}
\begin{itemize}[leftmargin=1.4em,itemsep=2pt]
\item 교환법칙: $A\cup B=B\cup A$, $A\cap B=B\cap A$
\item 결합법칙: $(A\cup B)\cup C=A\cup(B\cup C)$, $(A\cap B)\cap C=A\cap(B\cap C)$
\item 분배법칙: $A\cap(B\cup C)=(A\cap B)\cup(A\cap C)$, $A\cup(B\cap C)=(A\cup B)\cap(A\cup C)$
\end{itemize}
\textbf{\color{goldstrong}문제 4, 5, 6}은 벤 다이어그램을 색칠하여 양변이 같은 영역임을 확인하는 문제 --- 정답은 색칠된 벤 다이어그램(교사용 참고자료 별도 확인).
\end{answerbox}

\section{성취수준 (이 차시 범위)}
\begin{tabularx}{\linewidth}{@{}>{\bfseries\color{goldstrong}}p{0.9cm}X@{}}
\toprule
A & 교집합·합집합의 연산 및 원소의 개수 공식을 체계적으로 수행하고, 연산 법칙을 벤 다이어그램으로 확인할 수 있다. \\
B & 원소의 개수 공식을 활용하고, 연산 법칙을 안다. \\
C & 원소의 개수 공식을 활용할 수 있다. \\
D & 간단한 원소의 개수 공식을 활용할 수 있다. \\
E & 간단한 교집합·합집합 연산을 수행할 수 있다. \\
\bottomrule
\end{tabularx}

\section{다음 차시}
\begin{calloutbox}
\textbf{04 여집합과 차집합 --- 1/2차시.} `03 교집합과 합집합'이 여기서 마무리된다. 다음 시간부터는 전체집합을 기준으로 `제외'를 다루는 여집합과 차집합을 배운다.
\end{calloutbox}
\end{document}
